\section{Results}
\label{sec:results}

\subsection{Simple and Intervaled Results}
\label{ssec:res_ourmethod}
Explain the results of the simple and intervaled framework (cancer).
Compare simple vs intervaled. Get rate-distortion curves.
compare intervaled vs small update interval.

\figure[ht]
    \centering
%    \includegraphics[width=0.5\textwidth]{figures/results.png}
    \caption{Simple vs Intervaled Results}
    \label{fig:simple_vs_intervaled_rate_distortion}
    The figure shows the rate-distortion curves for the simple and intervaled frameworks.
\endfigure

\figure[ht]
    \centering
%    \includegraphics[width=0.5\textwidth]{figures/results.png}
    \caption{Theoretical vs Empirical Results}
    \label{fig:theoretical_vs_empirical_rate_distortion}
    The figure shows the rate-distortion curves for the theoretical and empirical results.
\endfigure

\figure[ht]
    \centering
%    \includegraphics[width=0.5\textwidth]{figures/results.png}
    \caption{compare with conventional quantizers}
    \label{fig:intervaled_vs_conventional_auc_rate_vs_rounds}
    The figure shows the AUC vs rounds and rate vs rounds for the intervaled framework and conventional quantizers.
\endfigure

\subsection{Performance Comparison}
\label{ssec:comparison}
Explain why no direct comparison due to no DSC quantizer experiments in literature.
But we can compare with the performance of the same model trained in a non-DSC quantizer setting.

\table[ht]
    \centering
    \begin{tabular}{|c|c|c|c|c|}
        \hline
        Method & AUC@R20 & AUC@R50 & Rate (bits) & MAPE \\
        \hline
        No Compression & 0.85 & 0.87 & N/A & N/A
    \end{tabular}
    \caption{Performance Comparison}
    \label{tab:performance_comparison}
    The table shows the performance comparison between our method and the no compression baseline.
\endtable